% Kamp 1968, Chapter IV, Theorems 4 and 5, printed pp. 108--112.

\subsection{Integer irredundancy and the rational-cut obstruction}
\label{sec:p09-kamp-iv4-iv5-separations}

\subsubsection{The converse integer irredundancy}

\begin{hbdefinition}{def:p09-kamp-iv4-basis}{Isolation and unbounded past/future without immediate neighbors}
On \(\mathsf{In}=\langle\mathbb Z,<\rangle\), retain the
unary operations \(\mathsf{ISOLP},\mathsf{ISOLF}\) with the
exact integer clauses in \cref{thm:p09-kamp-iv2-unary-semantics}.
For \(A\subseteq\mathbb Z\), define
\[
\begin{aligned}
 \mathsf{PAST}(A)&=\{i:\exists j<i\ (j\in A)\},\\
 \mathsf{FUTURE}(A)&=\{i:\exists j>i\ (j\in A)\}.
\end{aligned}
\]
These are strict existential operations without a bound on the
distance to the witness. The basis for Theorem IV.4 is exactly
\[
 \mathcal S''=
   \{\mathsf{ISOLP},\mathsf{ISOLF},\mathsf{PAST},
                       \mathsf{FUTURE},\mathsf{NOT},\mathsf{AND}\}.
\]
Its finite sentential formulas start with proposition letters and
use these four unary operations, negation, and conjunction.
The targets \(\mathsf{PAST}_1,\mathsf{FUTURE}_1\) inspect
the immediate integer neighbors, so
\(\mathsf{PAST}_1(A)=\{i:i-1\in A\}\) and
\(\mathsf{FUTURE}_1(A)=\{i:i+1\in A\}\).

For each integer \(r\), put
\(R_r=\{4k+r:k\in\mathbb Z\}\), so that
\(R_{r+4}=R_r\). Write
\(V=R_0\cup R_2\) for the even integers and
\(W=R_1\cup R_3\) for the odd integers.
Complements in this integer subsection are relative to \(\mathbb Z\)
\cite[ch.~IV, thm.~4, pp.~108--109]{Kamp1968}.
\end{hbdefinition}

\begin{hblemma}{thm:p09-kamp-iv4-field}{proved=annex:p9.kamp.iv4.field}{The complete eight-set calculation modulo four}
The family
\[
 \mathcal B_4=
 \{\mathbb Z,\varnothing,R_2,R_2^c,R_0,R_0^c,V,W\}
\]
is a field of subsets of \(\mathbb Z\). It is closed under
every operation of \(\mathcal S''\). The complete values of
the four unary operations on this family are
\[
\begin{array}{c|cc|cc}
 A&\mathsf{ISOLP}(A)&\mathsf{ISOLF}(A)
     &\mathsf{PAST}(A)&\mathsf{FUTURE}(A)\\ \hline
 \mathbb Z&\varnothing&\varnothing&\mathbb Z&\mathbb Z\\
 \varnothing&\varnothing&\varnothing&\varnothing&\varnothing\\
 R_2&\mathbb Z&\mathbb Z&\mathbb Z&\mathbb Z\\
 R_2^c&R_0&R_0&\mathbb Z&\mathbb Z\\
 R_0&\mathbb Z&\mathbb Z&\mathbb Z&\mathbb Z\\
 R_0^c&R_2&R_2&\mathbb Z&\mathbb Z\\
 V&\mathbb Z&\mathbb Z&\mathbb Z&\mathbb Z\\
 W&\mathbb Z&\mathbb Z&\mathbb Z&\mathbb Z
\end{array}
\]
The three atoms of this field are \(R_0,R_2,W\), so no
member distinguishes residue one from residue three. In particular,
if the only occurring input letter is interpreted by \(R_2\),
every finite \(\mathcal S''\)-formula has its truth set in
\(\mathcal B_4\). The same statement holds when every
occurring atomic interpretation already belongs to the field
\cite[ch.~IV, pp.~108--109]{Kamp1968}.
\end{hblemma}

\begin{hbtheorem}{thm:p09-kamp-iv4}{proved=annex:p9.kamp.iv4}{Immediate past and future are not definable from the remaining integer basis}
Neither \(\mathsf{PAST}_1\) nor \(\mathsf{FUTURE}_1\)
is expressible on \(\mathsf{In}\) by a finite one-input
formula over \(\mathcal S''\). Equivalence would have to
hold for every \(A\subseteq\mathbb Z\) at every integer.
The fixed input \(A=R_2\) already prevents either definition:
\[
              \mathsf{PAST}_1(R_2)=R_3,
       \qquad \mathsf{FUTURE}_1(R_2)=R_1,
\]
and neither target set belongs to \(\mathcal B_4\).
At the two current points one and three, every candidate formula
has equal truth values, whereas each target has unequal truth values
\cite[ch.~IV, thm.~4, pp.~108--109]{Kamp1968}.
\end{hbtheorem}

\subsubsection{Two interpretations on the rational order}

\begin{hbdefinition}{def:p09-kamp-rational-cut-models}{The periodic and accumulating rational propositions}
Now the fixed time domain is
\(\mathsf{Ra}=\langle\mathbb Q,<\rangle\).
For real numbers \(a<b\), write
\((a,b)_{\mathbb Q}=\{r\in\mathbb Q:a<r<b\}\).
Thus real interval endpoints specify subsets of the rational domain;
they are not extra time points or constants in the first-order language.
Put \(\alpha=\sqrt2\) and \(b_n=\alpha+1/n\) for
integers \(n\geq1\). All union indices below are integers.
Define
\[
\begin{aligned}
 Q&=\bigcup_{k\in\mathbb Z}
                   (\alpha+2k,\alpha+2k+1)_{\mathbb Q},\\
 Q'&=\bigcup_{\substack{k\in\mathbb Z\\k\ne0}}
                   (\alpha+2k,\alpha+2k+1)_{\mathbb Q}\\
   &\quad\cup\
       \bigcup_{j\geq0}(b_{2j+2},b_{2j+1})_{\mathbb Q}.
\end{aligned}
\]
The second interpretation replaces the original band
\((\alpha,\alpha+1)_{\mathbb Q}\) by alternating bands
whose endpoints decrease toward the irrational cut \(\alpha\).
In the displayed intervals the smaller endpoint is always first.

Let \(\mathcal M=\langle\mathbb Q,<,Q\rangle\) and
\(\mathcal M'=\langle\mathbb Q,<,Q'\rangle\), interpreting
the sole unary predicate \(P_0\) by these sets. The same sets
interpret \(q_0\) in the source's temporal language
\(\mathsf{TL}_1\), whose constructors are negation,
conjunction, strict \(S\), and strict \(U\), with the
witness operand first and the open-interval operand second.
Complements from this point onward are relative to \(\mathbb Q\)
\cite[ch.~IV, thm.~5 construction, pp.~109--110]{Kamp1968}.
\end{hbdefinition}

\begin{hblemma}{thm:p09-kamp-rational-local-table}{proved=annex:p9.kamp.rational.local.table}{All temporal operations on the four local truth-set choices}
Each of \(Q,Q'\) is locally constant on the rational domain:
every rational point has a rational open neighborhood on which its
membership is constant. Each set and its complement are unbounded
in both directions. Neither an irrational band endpoint nor the
accumulation cut is a point of the time domain.

More generally, let \(A\subseteq\mathbb Q\) have these
local-constancy and two-sided unboundedness properties. For
\(X,Y\in\{\varnothing,A,A^c,\mathbb Q\}\), both
strict temporal operations have the following complete table,
with the witness set \(X\) indexing rows and the interval
set \(Y\) indexing columns:
\[
\begin{array}{c|cccc}
 S(X,Y)=U(X,Y)&\varnothing&A&A^c&\mathbb Q\\ \hline
 \varnothing&\varnothing&\varnothing&\varnothing&\varnothing\\
 A&\varnothing&A&\varnothing&\mathbb Q\\
 A^c&\varnothing&\varnothing&A^c&\mathbb Q\\
 \mathbb Q&\varnothing&A&A^c&\mathbb Q
\end{array}
\]
The same table therefore governs the two specified interpretations
\cite[ch.~IV, pp.~110--112]{Kamp1968}.
\end{hblemma}

\begin{hblemma}{thm:p09-kamp-iv5-temporal-induction}{proved=annex:p9.kamp.iv5.temporal.induction}{The full one-input temporal indistinguishability induction}
For every finite \(\mathsf{TL}_1\)-formula \(\eta\) with
no proposition letters other than \(q_0\), the pair of truth sets
\[
  \bigl(\llbracket\eta\rrbracket_{\mathcal M},
        \llbracket\eta\rrbracket_{\mathcal M'}\bigr)
\]
is one of
\[
 (\varnothing,\varnothing),\quad
 (\mathbb Q,\mathbb Q),\quad (Q,Q'),\quad (Q^c,(Q')^c).
\]
Consequently all three source comparison assertions hold for any
\(r,s\in\mathbb Q\):

\begin{enumerate}
\item If \(r\in Q\) exactly when \(s\in Q\), then
      \(\eta\) has equal truth values at \(r,s\) in
      \(\mathcal M\).
\item If \(r\in Q'\) exactly when \(s\in Q'\), then
      \(\eta\) has equal truth values at \(r,s\) in
      \(\mathcal M'\).
\item If \(r\in Q\) exactly when \(s\in Q'\), then
      \(\eta\)'s truth at \(r\) in \(\mathcal M\)
      equals its truth at \(s\) in \(\mathcal M'\).
\end{enumerate}

In particular every such formula agrees at the rational current point
three in the two interpretations
\cite[ch.~IV, thm.~5 proof, pp.~110--112]{Kamp1968}.
\end{hblemma}

\begin{hbdefinition}{def:p09-kamp-rational-distinguishing-formula}{A first-order condition for the accumulating obstruction}
The following equality-free formula uses only \(<,P_0\), and
has exactly the current point \(t_0\) free:
\[
\begin{aligned}
 \beta(t_0)=\exists t_1\bigl[&t_1<t_0\land\neg P_0(t_1)\land
           \exists t_2(t_1<t_2<t_0\land P_0(t_2))\\
 &\land\forall t_2\bigl((t_1<t_2<t_0\land P_0(t_2))\to\\
 &\hspace{18mm}\exists t_3\exists t_4
      (t_1<t_3<t_4<t_2\land P_0(t_3)\land\neg P_0(t_4))\bigr)
 \bigr].
\end{aligned}
\]
All quantified variables range over \(\mathbb Q\). Put
\(\varphi(t_0)=\neg\beta(t_0)\).
The witness \(t_1\) for \(\beta\) is a false point with
a true point between it and the current time. Every true point
in that interval must have, still earlier in the interval, a
true point followed by a false point. Thus no first true block
after \(t_1\) can account for all subsequent true points.
The existential true-point clause excludes a vacuous instance.

This gives a corrected distinguishing formula for the source's
two-interpretation construction. Its definition contains no
\(\sqrt2\), real endpoint, or named rational parameter;
those enter only the chosen interpretations used to prove separation.
\end{hbdefinition}

\begin{hbtheorem}{thm:p09-kamp-iv5}{proved=annex:p9.kamp.iv5}{A first-order unary rational tense beyond SINCE and UNTIL}
The formula \(\varphi\) defines a unary
\(\mathsf{Ra}\)-tense which is not expressible in terms of
\(\mathsf{NOT},\mathsf{AND},\mathsf{SINCE},\mathsf{UNTIL}\).
Precisely, there is no finite one-input \(\mathsf{TL}_1\)
formula equivalent to \(\varphi\) for every unary predicate
interpretation on \(\mathbb Q\) and every rational current point.
For the explicit interpretations above,
\[
                     \mathcal M\models\varphi[t_0\mapsto3],
            \qquad \mathcal M'\not\models\varphi[t_0\mapsto3],
\]
while every one-input temporal formula agrees there. The order
domain remains \(\mathbb Q\) throughout; the real cut is not
adjoined as a possible time point. Hence the complete-order
expressive-completeness theorem cannot be extended unchanged to
all linear orders
\cite[ch.~IV, thm.~5, pp.~109--112]{Kamp1968}.
\end{hbtheorem}

\begin{hbexample}{def:p09-kamp-iv5-cut-example}{Concrete witnesses around the missing point}
Since \(1<\sqrt2<3/2\), the rational point one lies in
the false band immediately below \(\alpha\), and two lies
in the true band \((b_2,b_1)_{\mathbb Q}\) in both models.
The current point three lies in
\((\alpha+1,\alpha+2)_{\mathbb Q}\), which is false in
both. In \(\mathcal M'\), take \(t_1=1\) as the witness
for \(\beta(3)\). Before every true point between one and
three, sufficiently small reciprocal bands supply a true point
followed by a false point. In \(\mathcal M\), any false
starting point with a later true point has a first subsequent
true band; a point inside that band prevents the universal clause
of \(\beta\). The annex proves these assertions for every
quantified point, including all rational positions near the cut.
\end{hbexample}
